\section{Motivation}

\begin{frame}
    \frametitle{Microservice vs Monolith}

    \begin{tikzpicture}

        \node (test) [circle, fill=green] {};

    \end{tikzpicture}
    Bild Monolioth vs Microservice
    
\end{frame}

\begin{frame}
    \frametitle{Herausforderungen von Microsverices}

    Microservices bringen Nachteile von verteilten Systemen mit sich. Durch die Kommunikation übers Netzwerk insbesondere:

    \begin{itemize}
        
        \item Microservices können ausfallen, nicht erreichbar sein
        \item Aufrufe können verloren gehen
        \item Aufrufe können eine hohe Latenz aufweisen
    \end{itemize}
\end{frame}

\begin{frame}
    \frametitle{Entuwrfsmuster zur Erhöhung der Resilienz}

    \begin{columns}
        \begin{column}{.3\textwidth}
            Entwurfsmuster zum Umgang mit benannten Fehlerquellen:

            \begin{itemize}
                \item Retry
                \item Timeout
                \item Fallback
                \item Circuit Breaker
            \end{itemize}
        \end{column}

        \begin{column}{.7\textwidth}
            \begin{tikzpicture}[align=center, node distance=3cm and 2cm]

                \node (closed) [state] {geschlossen};
                \node (open) [state, right=of closed] {offen};
                \node (halfopen) [state, below=of open] {halboffen};
                
                \draw [arrow] (closed) -- node[anchor=south, text width=3cm, yshift=5mm] {Schwellwert überschritten} (open);
                \draw [arrow] ([xshift=-10mm]open.south) -- node[anchor=east, text width=3cm, xshift=5mm] {Reset Zeit abgelaufen} ([xshift=-10mm]halfopen.north);
                \draw [arrow] ([xshift=10mm]halfopen.north) -- node[anchor=west, text width=3cm, xshift=-5mm] {Aufruf war fehlerhaft} ([xshift=10mm]open.south);
                
                \draw [arrow] (halfopen) -| node[anchor=north] {Aufruf war erfolgreich} (closed);
            
            \end{tikzpicture}
        \end{column}
    \end{columns}

\end{frame}


\begin{frame}
    \begin{large}
    \large Wie können die Muster effektiv eingesetzt und entsprechend konfiguriert werden?

    \end{large}
\end{frame}